\documentclass{article}
\PassOptionsToPackage{numbers,compress}{natbib}
\usepackage[preprint]{neurips_2025}
\usepackage{booktabs}
\usepackage{graphicx}
\usepackage{hyperref}
\usepackage{microtype}
\usepackage{amsmath}

\title{Evolving Vision--Language Inference for Scientific Chart Understanding}
\author{Anonymous challenge submission}
\date{}

\begin{document}
\maketitle

\begin{abstract}
We study code-level inference optimization for scientific-chart question answering
with the fixed Qwen3-VL-2B-Instruct and Qwen3-VL-2B-Thinking checkpoints.
The system combines two manually designed programs with an AlphaEvolve-style,
catalog-constrained program search. Candidate programs are declarative, rendered
by a trusted template, evaluated in restricted subprocesses, and selected by
accuracy and measured latency. Development used CharXiv descriptive questions,
greedy decoding, and one physical H200 GPU. A first formal campaign exposed a
censored Thinking-baseline timeout; we retained it as a negative result and
pre-registered a Recovery-v2 campaign with per-query atomic journals. External
GPU contention later invalidated the whole first Speed32 block. We retained its
valid single-query and Dev128 records, then reran all 15 Speed32 program--repeat
blocks under an authorized, uncontended GPU0-only replacement. Every final number
below is generated from that sealed composition; no old and new speed records are
mixed. The held-out challenge test labels remain unavailable, so we make no
hidden-test claim. A subsequent portability-only repair removed server-bound
resource assumptions and replayed raw evidence independently. Earlier relocated
single-query checks reproduced the sealed outputs, but later executable/package
changes formally superseded those receipts. They remain historical
\textsc{portability-smoke-only} evidence and do not enter any metric or selection.
\end{abstract}

\section{Introduction}
Vision--language models (VLMs) can answer questions about scientific figures, but
their inference programs expose consequential choices beyond model weights:
message construction, image preprocessing, token budgets, and answer extraction.
Those choices affect both exact-match accuracy and wall-clock cost. This challenge
asks for manual and automatically evolved inference programs for two fixed Qwen3-VL
checkpoints. Our objective is therefore not training; it is controlled program
search under an immutable model, a fixed metric, and reproducible hardware
conditions.

Our main engineering contribution is an evidence-oriented evolution loop inspired
by AlphaEvolve~\citep{alphaevolve2025}. A local Qwen model proposes one option
from a pre-registered catalog, while a typed declarative specification and trusted
renderer prevent arbitrary generated Python from executing. Population state,
lineage, rendered source diffs, Pareto selection, checkpoints, and resume behavior
are explicit artifacts. Our empirical contribution separates Phase~2 search
evidence from a formal Recovery-v2 evaluation and its whole-set Speed32
replacement. This separation matters:
selection-set gains are useful for choosing candidates, but are not substitutes
for evaluation on the frozen 128-query development set.

\paragraph{Contributions.}
We provide (i) naive, manual, and evolved inference programs for both fixed
checkpoints; (ii) a restricted, catalog-constrained two-generation search with
replay-safe checkpoints; (iii) a formal protocol that measures accuracy once and
speed three times with CUDA-boundary timing; and (iv) a hash-linked reproducibility
and packaging path whose checks fail closed if formal evidence is absent or
modified.

\section{Task, Data, and Integrity Constraints}
CharXiv contains realistic figures and questions that probe chart
understanding~\citep{wang2024charxiv}. We use the descriptive validation JSON in
its original insertion order. Exact match is computed after stripping leading
and trailing whitespace and applying case-insensitive comparison; exceptions count
as errors and remain in the denominator. Development and final manifests are
frozen by query key and image hash. The challenge's hidden test data and labels
are neither present nor inferred.

The model IDs and revisions are fixed for Qwen3-VL-2B-Instruct and
Qwen3-VL-2B-Thinking~\citep{bai2025qwen3vl}. All evaluated programs use BF16,
SDPA, batch size one, and greedy generation
(\texttt{do\_sample=false}, \texttt{num\_beams=1}, and one returned sequence).
No weights are trained or altered. Locked snapshot manifests bind the model
configuration, processor, tokenizer, chat template, and weight shards; model and
processor resolve from the same source and revision. BF16 is the execution dtype,
not a quantized or modified checkpoint. All \emph{formal} GPU experiments run
sequentially with \texttt{CUDA\_VISIBLE\_DEVICES=0}, one visible device, and
logical \texttt{cuda:0}; GPU~1 is excluded from formal measurement and selection.
After all formal results were frozen, an explicitly non-formal portability check
used physical GPU~0 for the two Instruct programs and physical GPU~1 for the two
Thinking programs. Each lane exposed one logical \texttt{cuda:0} and ran its two
programs serially. Model loading is excluded from steady-state formal timing.

\section{Inference Programs}
\subsection{Naive baselines}
The naive programs preserve the challenge's direct model/processor structure and
serve as reference points for both correctness and speed. They use the checkpoint
chat template, one image, one question, greedy decoding, and a conservative answer
normalization step. They intentionally avoid search-derived prompt or extraction
choices. The Instruct baseline is the accuracy floor used by two of the final
selection rules. The Thinking baseline remains part of the six-program single
smoke evaluation, but its earlier Dev128 run exceeded 7200 seconds and is explicitly
reported as \textsc{timeout-censored}; no accuracy is imputed.

\subsection{Manual programs}
The manual Instruct program requests a concise, explicitly tagged final answer and
uses a conservative tag-aware parser with a 48-token cap. The manual Thinking
program preserves native reasoning behavior, requests a tagged final answer,
stops on the closing tag when available, and uses the single pre-registered
384-token arm. That arm followed an initial four-query result in which longer
reasoning frequently failed to yield a parseable final answer. The arm changed
only the token cap; it did not use question IDs, labels, or sample-specific rules.

On the Phase~2 four-query smoke set, manual Instruct answered 3/4 and manual
Thinking answered 0/4. These small results were operational gates, not model
selection claims. They are included to document that both programs performed
real inference and to make the negative Thinking outcome visible. Final
comparisons use the formal table generated below.

\section{Catalog-Constrained Evolution}
\subsection{Why declarative candidates}
Early provider attempts established an important negative result. Free-form local
mutation could return a duplicate or an off-focus edit even after one correction.
Those completions were retained as failures and were never re-labeled as children.
The production search therefore restricts each proposal to a typed
\texttt{CANDIDATE} literal. Fields cover prompt template, system instruction,
reasoning mode, token budget, preprocessing, and answer extraction, with checkpoint-
specific ranges. Model ID, BF16, SDPA, batch size one, and greedy parameters are
hard-locked. Paths, imports, calls, attributes, shell commands, network access,
query identifiers, and answers are not candidate dimensions.

A trusted renderer maps a validated specification to complete, deterministic
\texttt{vlm\_inference(image\_path, question)} source. Lineage and diffs are
computed from that rendered source. Semantic canonicalization removes dimensions
that would not change runtime behavior. A proposal is valid only when the
pre-registered focus changes, all non-focus fields remain permitted, the rendered
source passes static/import checks, and restricted single-query evaluation
succeeds.

\subsection{Local catalog provider}
For each request, a fixed local Instruct checkpoint scores a frozen ordered option
catalog. The provider uses the exact chat suffix and next-token logits for
single-token labels. Before GPU use, labels are verified to be distinct one-token
encodings in the exact context. Scores are preserved as float32 bytes, hexadecimal
floats, and a full ordinal ranking; ties break only by catalog ordinal. A runtime-
invalid top option fails that slot. Only a duplicate already present in the parent
or population may advance to the next option in the same stored ranking, without
another model call. Network and credential-bearing environment variables are
removed, and provider processes exit before target evaluation begins.

Qualification tries a fixed focus order and stops at the first valid receipt.
Generation~1 creates two children from the manual seed. Generation~2 creates up
to four children from at most two evaluated parents, with at least one valid
generation-1 child represented when available. Each model is limited to six
children and two generations. No third generation, hidden-test query, or
post-hoc search arm is allowed.

\subsection{Restricted evaluation and population state}
The parent process alone owns ground truth and scoring. A target worker receives
an opaque request ID, question text, and image bytes---not a split name, query key,
path, expected answer, or correctness bit. Candidate code is never executed
directly: only the trusted renderer/runtime is invoked. The subprocess runs with
a minimized environment, Evo-local working/cache/temp directories,
\texttt{shell=false}, a new process group, bounded frames, finite-number checks,
and a timeout that kills only its own group. Tests verify that an unrelated
sentinel survives timeout handling.

The population tracks parent IDs, generation, proposal request IDs, catalog scores,
specifications, source hashes, source diffs, single-query gates, Search12 metrics,
and failures. Parent selection is multi-objective: valid candidates enter a
Pareto comparison of correctness and latency. Generation~1 is atomically
checkpointed; a fresh resume creates a pre-GPU generation~2 journal. Once marked
complete, a second resume returns before any provider or GPU call. Crash-injection
tests cover checkpoint, export, summary, and completed-marker boundaries.

\subsection{Phase 2 search evidence}
For Instruct, two generation-1 and four generation-2 proposals were valid. On the
fixed Search12 set, scores ranged from 5/12 to 7/12, while the manual seed scored
6/12. Selection32 evaluated the seed and two real children under one identical
manifest and warmup. All three scored 17/32; the exported real child had a mean
latency 0.024920 seconds lower than the seed on that development selection set.
For Thinking, six children were valid; Search12 scores ranged from 2/12 to 3/12
against a 3/12 seed. On Selection32, the seed scored 7/32 and both finalists
scored 12/32. The exported child was 0.017065 seconds slower in mean latency than
the seed. These are selection facts, not statistical or held-out claims.

The exact evolved exports are immutable rendered children. We retain the manual
seed in the same Selection32 protocol to prevent a child-only advantage. Provider
and target network-attempt counts were zero, all target errors in the cited Phase~2
selection runs were zero, and completed resume checks left canonical artifact
hashes and mtimes unchanged.

\section{Formal Recovery-v2 Evaluation}
\subsection{Pre-registration and manifests}
Recovery-v2 was frozen before inspecting its results. The accuracy manifest is
the first 128 descriptive validation queries in source order. Evolution-used
queries are recorded separately, and an untouched complement is reported without
being used for reselection. Speed32 is deterministically hash-selected, with at
most one query per figure where possible; its warmup figure is disjoint. Programs
use identical query order, image bytes, dtype, generation semantics, and error
handling.

The formal candidate set contains six programs for single-query smoke:
two naive, two manual, and two evolved. Dev128 and Speed32 include naive Instruct
plus the four optimized programs. Naive Thinking is not silently dropped: its
v1 Dev128 event is retained as \textsc{timeout-censored} above 7200 seconds with
accuracy unavailable. Recovery-v2 gives each scheduled program a 14400-second
wall bound and continues to the next program after a terminal timeout.

\subsection{Measurement}
Each Dev128 program is evaluated exactly once. Speed32 uses three complete,
steady-state repeats in a frozen rotating program order; no fastest repeat is
selected. For every query, a non-scored warmup precedes measurement, CUDA is
synchronized immediately before and after the timed region, and model/processor
load is excluded. The journal records preprocessing, generation, and total
seconds, generated-token count, termination reason, peak memory, response hash,
and error state. Speed selection first takes the median of three repeats for each
query, then aggregates those 32 per-query medians into median, mean, Type-7 p95,
and queries per second.

Each query journal record is committed outside the measured interval with atomic
no-replace semantics. Resume accepts only a contiguous verified prefix whose
source, program, manifest, and row hashes match the pre-registered run. A partial
query is not guessed or overwritten. Synthetic tests cover interruption after
several queries, duplicate keys, reordered rows, numeric tampering, source
tampering, and finalization after all 26 scheduled results.

\subsection{Whole-set Speed32 replacement}
After all 26 original Recovery-v2 runs completed, process telemetry showed that an
unrelated process overlapped GPU~0 during part of the speed campaign. We therefore
marked the \emph{entire} original Speed32 block non-qualifying; none of its 15
speed measurements enters selection. The 11 valid single-query and Dev128 records
remain immutable. An authorized replacement reran all five programs and all three
frozen repeats on hopper25 physical GPU~0, preserving the exact program sources,
Speed32 manifest, warmup, rotation, evaluator, and selection rules. A 120-second
idle gate preceded the run, and five-second telemetry observed no foreign GPU~0
process throughout the replacement. Validation requires exactly 11 retained
records plus 15 replacement records and rejects partial replacement, copied old
speed records, or any old/new mix.

\subsection{Results}
The following table is generated by
\texttt{evaluation/recovery\_finalize.py} from the sealed Recovery-v2 composition.
Compilation intentionally fails if that generated evidence is absent.

% Generated from the semantically verified Recovery-v2 composition; do not edit.
\begin{table}[t]
\centering\small
\caption{Recovery-v2 formal validation with authorized whole-set Speed32 replacement. Accuracy is exact match on Dev128; latency is the median of the 32 per-query medians over three steady-state repeats.}
\label{tab:formal-results}
\begin{tabular}{lrrrr}
\toprule
Program & Correct/128 & Errors & Median (s) & Type-7 p95 (s) \\
\midrule
baseline\_instruct & 51/128 & 0 & 14.541101 & 19.728077 \\
manual\_instruct & 78/128 & 0 & 15.443962 & 21.594249 \\
manual\_thinking & 23/128 & 0 & 20.043467 & 27.825051 \\
evolved\_instruct & 78/128 & 0 & 14.071758 & 19.757527 \\
evolved\_thinking & 44/128 & 0 & 17.912906 & 26.498611 \\
\bottomrule
\end{tabular}
\end{table}

\begin{table}[t]
\centering\small
\caption{Pre-registered optimized-program selections.}
\label{tab:best-programs}
\begin{tabular}{ll}
\toprule
Role & Selected program \\
\midrule
best\_accuracy & evolved\_instruct \\
best\_speed & evolved\_instruct \\
best\_overall & evolved\_instruct \\
\bottomrule
\end{tabular}
\end{table}


Accuracy is reported as an integer correct count over 128 as well as a fraction.
All speed repeats are retained. Paired comparisons, a fixed-seed bootstrap, and
McNemar counts are computed from aligned per-query records; the complement analysis
is descriptive and never changes the selected programs.

The result pattern also identifies which bundled techniques are associated with
the final roles, without pretending that this study is a field-by-field causal
ablation. Relative to naive Instruct, both optimized Instruct programs gain 27
exact matches; their principal shared intervention is a concise, tagged-answer
prompt. Evolution retained that prompt while changing the manual 48-token,
tag-aware path to a 64-token budget with boxed-or-final extraction. The evolved
program preserves 78/128 and has the lowest qualifying median latency
(14.071758~s), versus 15.443962~s for manual Instruct. For Thinking, the evolved
visible-evidence prompt and terminal-punctuation postprocessing accompany a gain
from 23/128 to 44/128 and a lower median (17.912906~s versus 20.043467~s), but that
program remains below the Instruct accuracy leaders. Because several choices
change together, these comparisons support bundle-level attribution only.

\subsection{Frozen best-program rules}
Only the four optimized programs are eligible for the three \texttt{best\_*}
exports. \textbf{Best accuracy} maximizes integer correct count, then minimizes
errors, speed median, generated tokens, and source hash. \textbf{Best speed}
requires zero errors and accuracy at least
$\max(\text{naive-Instruct correct},\text{best-accuracy correct}-4)$; among eligible
programs it minimizes median speed, p95, tokens, then maximizes accuracy and uses
the source hash. Four correct answers equal $1/32$ on the 128-query accuracy set,
matching the frozen tolerance.

\textbf{Best overall} is author-defined because the challenge specifies no official
joint score. Eligibility requires zero errors and accuracy no worse than naive
Instruct. For each eligible program,
$A=\mathrm{accuracy}/\max \mathrm{accuracy}$ and
$S=\min \mathrm{latency}/\mathrm{latency}$; the score is their harmonic mean.
Ties prefer accuracy, then latency, tokens, and source hash. Finalization
recomputes all three rules, copies the selected source byte-for-byte, and seals
the table, numeric trace, campaign evidence, source identity, and exports in one
manifest.

All three rules select \texttt{evolved\_instruct}; this is one genuine winner
serving three roles, not three independently claimed variants. It wins accuracy
by the frozen tie-break over the equally accurate manual Instruct program, wins
speed because its median is the lowest among accuracy-eligible programs, and wins
overall because it combines normalized maximum accuracy with normalized maximum
eligible speed. Consequently, the concise tagged prompt plus robust extraction
most directly support the accuracy role, the measured low-latency implementation
supports the speed role, and their conjunction supports the overall role.

\section{Negative Results and Failure Analysis}
Three negative classes shaped the final design. First, unrestricted local
proposal generation produced duplicate and off-focus specifications. The system
consumed those requests, saved their hashes and reason codes, and did not fabricate
children. Catalog scoring made the choice space auditable and bounded. Second,
launcher and isolated-bytecode failures occurred before measurement in development.
They were retained, verified to have no measurement artifacts, and did not count
as successful proposals. Resume then used the same frozen campaign rather than
silently restarting it.

Third, formal v1 completed seven results before naive Thinking Dev128 remained
healthy but exceeded the 7200-second outer limit. The runner terminated it,
recorded no complete raw/summary/journal, released GPU~0, and did not auto-repeat.
That observation is right-censored: it establishes only total time greater than
7200 seconds and supplies no accuracy estimate. Recovery-v2 preserves the entire
v1 directory byte-for-byte, gives each new program a larger explicit bound, adds
per-query journals, and continues other scheduled programs after timeout.
The immutable v1 closeout and Recovery-v2 protocol had declared the wrong v1
campaign-plan digest; the verified erratum binds their unchanged bytes to the
actual plan SHA-256 \texttt{effe1b6a1741885b2ab4c48c71efeedc1735825039d837ea0533b79f13258c79},
with neither formal metrics nor selection affected.

Finally, a foreign GPU~0 process overlapped the original Recovery-v2 speed window.
We preserved all 26 original artifacts byte-for-byte and recorded an append-only
contamination incident rather than selecting from questionable latency data. The
authorized whole-set replacement was accepted only after all 15 blocks completed
under the frozen protocol with uncontended telemetry. Dev128 accuracy was retained;
the original Speed32 block remains permanently non-qualifying.

\section{Security, Validity, and Limitations}
The restricted worker meaningfully reduces accidental leakage and code-execution
risk, but it is not a kernel security boundary. Namespace capability is probed
and reported; where unavailable, only trusted catalog-rendered candidates are
eligible. The provider is offline and local, candidate inputs omit labels and
identifiers, and all cache/temp/history paths remain inside the project root.
Symlink audits resolve every link and reject absolute, broken, external, escape,
or looping targets. Tests also exercise malformed frames, duplicate JSON keys,
NaN/Infinity, credential/proxy stripping, target timeouts, unrelated-process
survival, artifact races, and no-overwrite publication.

Internal validity is limited by one formal-measurement GPU, one dtype, batch size
one, and a modest development budget. Phase~2 Search12 and Selection32 were used for selection and
therefore are optimistic estimates. Dev128 is public validation data rather than
the hidden challenge test. It is also not a strictly untouched holdout: exact
queries \texttt{21\_2} and \texttt{68\_2} occur in the recorded evolution-used
set, and eight Dev128 queries share those two figures. Exact match ignores semantic equivalence and may
penalize explanatory yet correct responses. Latency on a shared cluster can vary
with system load even under synchronization and repeated measurement. We report
raw repeats and paired records so readers can inspect that variability.

External validity is also limited: both checkpoints are 2B Qwen3-VL variants and
all questions are CharXiv descriptive queries. Improvements may not transfer to
reasoning questions, other VLMs, quantized inference, or batched serving. We do
not estimate hidden-test accuracy, claim significance from Phase~2 selection
sets, or reinterpret timeout censoring as a measured score.

\section{Reproducibility and Artifact Traceability}
The repository contains locked dependencies, fixed model revisions and snapshot
manifests, data/image manifests, candidate sources, evolution lineage, tests,
formal protocols, and a single \texttt{reproduce.sh}. With no mode argument,
\texttt{reproduce.sh} selects \texttt{fresh} and invokes the official
\texttt{evaluate.py PROGRAM} entry in isolated GPU lanes. The explicit
\texttt{evidence} mode is a standard-library, CPU-only replay that reopens raw records,
recomputes correctness, errors, latency aggregates, speed selection, and all
three best-program rules, and refuses missing or changed evidence. Fresh inference
is therefore the operational default, while evidence replay remains independently
available without model loading. The final checker additionally requires the complete
Recovery-v2 evidence and authorized whole-set Speed32 replacement, verifies the sealed
11-plus-15 composition, checks the official NeurIPS PDF, and audits the deterministic ZIP.
Historical Phase~2 and v1 evidence have separate CPU-only verification modes and
cannot satisfy final completion.

\paragraph{Portability closeout.}
Package resources and managed outputs are physically contained, while explicit
model snapshots and evaluator images are validated as read-only regular files
instead of being forced beneath the source tree. Model and processor always
resolve from the same locked source; an invalid override fails without fallback.
The superseded root \texttt{BLOCKED\_STATUS.json} is retained only in a historical
archive and excluded from the current package. Two historical source snapshots
were extracted with \texttt{git archive} into fresh staging directories, where
four Single1 runs reproduced the sealed query, question, image hash, prediction
bytes, and correctness with zero errors. Subsequent executable/package changes
superseded both receipts as final gates, so they are not presented as proof for
the current archive. Current-package portability must instead pass clean-extraction
CPU imports, evidence replay, manifest/ZIP verification, and packaged-image path
checks; fresh GPU inference remains conditional on the two locked model snapshots
and suitable CUDA hardware. None of these checks changes formal accuracy or speed
claims.

\paragraph{Final fresh reproduction.}
After the portability source and evidence gates were frozen, one no-overwrite
reproduction ran five Dev128 programs and four Single1 checks through the official
entry point. The Dev128 manifest contained 128 ordered queries backed by 128
distinct regular-file paths and inodes. Those paths intentionally duplicate figure
bytes when different questions reference the same figure: the manifest discloses
32 unique image-content hashes rather than claiming 128 unique figures. The run
used the same locked model snapshots and unchanged weight bytes described above.
Its strict receipt closed all nine start/result/end chains, and a copied 42-file
evidence set passed independent receipt verification.

This fresh run is labeled \textsc{final-fresh-reproduction-only}. It changes no
formal Recovery-v2 metric or best-program selection. In particular, the fresh
Single1 check for naive Thinking does not replace or impute the censored
Dev128 Thinking baseline. The CPU-only \texttt{evidence} mode validates this fresh
receipt together with the retained formal evidence.

Every formal number in this paper is traceable through
\texttt{report/generated/report\_number\_trace.json}; every row names the source
artifact and hash. The final manifest binds both campaign IDs, the contamination
audit and replacement authorization, source commit/tree, old and replacement
plans, summaries and completed markers, all 26 composed result artifacts, official
style hash, table, trace, and byte-identical best exports. The checker reopens and rehashes those
files, recomputes metrics and selection, rejects fallback document classes or
missing formal evidence, and audits ZIP members for traversal, symlinks,
duplicates, oversize payloads, empty required files, and manifest mismatch.

\bibliographystyle{plainnat}
\bibliography{references}
\end{document}
